% ================= C 卷 =================
% 难度定位：拔高挑战（综合性强、区分度高）
\juantou{C}
\defenBiaoFour{20}{32}{30}{18}

\dati{名词解释}{每小题 5 分，共 20 分}

\begin{ti}
\xiaoti{1.（5 分）规律}
\kongbai{2.6cm}
\end{ti}

\begin{ti}
\xiaoti{2.（5 分）价值（哲学范畴）}
\kongbai{2.6cm}
\end{ti}

\begin{ti}
\xiaoti{3.（5 分）资本积累}
\kongbai{2.6cm}
\end{ti}

\begin{ti}
\xiaoti{4.（5 分）共产主义}
\kongbai{2.6cm}
\end{ti}

\dati{简答题}{每小题 8 分，共 32 分}

\begin{ti}
\xiaoti{1.（8 分）意识的能动作用表现在哪些方面？正确发挥这一作用需要具备什么条件？}
\kongbai{4.4cm}
\end{ti}

\begin{ti}
\xiaoti{2.（8 分）辩证否定观的基本内容是什么？它与形而上学否定观有何区别？}
\kongbai{4.4cm}
\end{ti}

\begin{ti}
\xiaoti{3.（8 分）剩余价值是怎样生产出来的？不变资本与可变资本在其中的作用有何不同？}
\kongbai{4.4cm}
\end{ti}

\begin{ti}
\xiaoti{4.（8 分）经济全球化具有二重性，请分别说明其积极作用与负面效应。}
\kongbai{4.4cm}
\end{ti}

\dati{论述题}{每小题 15 分，共 30 分}

\begin{ti}
\xiaoti{1.（15 分）为什么说社会基本矛盾是社会发展的根本动力？请结合生产力与生产关系、经济基础与上层建筑的矛盾运动加以论述。}
\kongbai{9.5cm}
\end{ti}

\begin{ti}
\xiaoti{2.（15 分）论述资本主义经济危机的实质与根源，并说明为什么生产社会化与生产资料私人占有的矛盾无法在资本主义制度内部根本解决。}
\kongbai{9.5cm}
\end{ti}

\dati{材料分析题}{共 18 分，阅读所给材料后回答问题}

\begin{ti}
\xiaoti{1.（18 分）马克思、恩格斯在《共产党宣言》中宣告资产阶级的灭亡和无产阶级的胜利是同样不可避免的；而后在《〈政治经济学批判〉序言》中指出，无论哪一个社会形态，在它所能容纳的全部生产力发挥出来以前，是决不会灭亡的。习近平总书记多次强调，我们依然处在马克思主义所指明的历史时代，既要树立共产主义远大理想，又要脚踏实地推进中国特色社会主义。\\

请回答：\\
（1）如何理解“两个必然”与“两个决不会”的内在统一？（9 分）\\
（2）这一思想对认识社会主义代替资本主义的长期性有何启示？（9 分）}
\kongbai{12.0cm}
\end{ti}
